\section{怎样翻译文言文}

把文言文译成准确、通顺而又生动的语体文，不仅需要掌握古代汉语的词义变化、各类句式和古代文化常识，而且要有较强的现代汉语的表达能力。所以，这是培养文言文阅读能力和语体文表达能力的一个好方法。

对于初学古文的中学生和社会青年来说，翻译的基本要求就是“准确”二字。能准确地理解原文的字、词、句在文中的含义，并能较快地读懂全文和用现代汉语准确地表述出来，这就说明已经具有翻译文言文的能力。如果译文是准确的，一般自会通顺的，至于流畅和生动，甚至能体现出原文的艺术特色和个人风格，则是以后逐步提高的事。

为了使译文能够准确地表达原文的意思，我们该怎么办呢？

做好翻译准备工作。
有的人拿到一篇文言文，就一字一句翻译起来，这是一种不好的习惯。因为脱离了前后贯串的语言环境，就不可能准确地理解词义和句法，结果往往是译不下去，或者译得歪曲了原意。例如，高考试卷中曾经有过这样一道试题，叫考生将《史记·淮阴侯列传》一段话中的加了横线的句子翻译成为现代汉语：

上（指汉高祖刘邦）尝从容与信（指韩信）言诸将能，各有差。上问曰：“如我，能将几何（ ）？”信曰：“陛下不过能将十万。”

上曰：“子公何如（ ）？”

曰：“如臣，多多而益善耳。”上笑曰：“多多益善，何为为我禽（ ）！”

信曰：“陛下不能将兵，而善将将（ ），此乃信之所以为陛下禽也（ ）。”

如果我们在动笔翻译之前，先通观全文，理清脉络，粗明大意，做到了成竹在胸；然后，根据上下文意和句子结构，弄清某些在文中有特殊用法的词语，突破难点；再分析文中在词序、句式等方面跟现代汉语的不同之处；做好了

这些融会贯通的准备工作之后，再着手翻译，那么，这几句文言句是并不难译好的。但是，不少考生只着眼于一个字一个句地孤立理解，而不管上下文意如何，就把“如我，能将几何？”翻译成“象我这样，有本能的将军有多少呢？”并把“于公何如？”翻译成“于将军（或“于先生”“于公”）怎么样？”这就不对了。因为“如我”与下文的“如臣”是对着说的，而“如臣，多多而益善耳”是回答“于公何如”的，并且是与“陛下不过能将十万”对着说的，如果再联系下边的“陛下不能将兵，而善将将”，就可以确定“如我，能将几何”中的“将”字是动词“带领、统率”的意思，句意是“象我，能带多少兵呢？”只有这样的翻译，才能与下文的“如臣，多多而益善耳”（象我，带得越多就越好啊）丝丝入扣，对举吻合。理解了这一点，“于公何如”这个与“如我，能将几何”对照而问的句子，只能翻译成“对你来说怎么样呢？”只有这样译，才能跟“如臣，多多而益善耳”的回答密切相合，又怎么可能跑出一个莫名其妙的什么“于将军”来呢？因为“公”字只是对于对方的尊称嘛。在贯串上下文意、理清前后脉络、突破难解字词的基础上，“陛下不能将兵，而善将将”当然就是“陛下您不擅长带兵，却善于统率将领”的意思了，这儿的关键在于理解“将兵”、“将将”的“将”字是动词“带领、统率”的意思。

剩下的“多多益善，何为为我禽”和“此乃信之所以为陛下禽也”两句话，前者是承“既然你带兵越多越好”而反问“何乃为我禽”，后者是承“陛下不能将兵，而善将将”而申述“这就是我韩信‘为陛下禽’的原因”。不少考生将它们翻译成“为什么成为我的家禽呢？”和“这就是我韩信

成为陛下家禽的原因。”这种不顾上下文意的译法恐怕连自己也莫名其妙而令人啼笑皆非了。既然韩信是刘邦手下的大将，又自称带兵越多越好，而说刘邦不过能带十万兵，那么刘邦的反问当然还是从带兵打仗的角度说的，绝不会说到什么“家禽”问题上去。因此，对于这个难以讲通的“禽”字，就要想到通假字上去。稍加回想，就会知道，原来这个“禽”字正是“将军禽操，宜在今日”（《赤壁之战》）、“阴知奸党名姓，一时收禽”（《张衡传》）等文章中常见通假字“擒”字。而“为我禽”、“为陛下禽”的“为”当然是作为被动词讲了。所以“为我禽”是“被我捉住”的意思，“为陛下禽”则是“被您捉住”的意思了。

俗话说：“磨刀不误砍柴工。”我们只有先贯通上下文意，做到了胸有全局，然后再推敲难字、词序和句式，把这些准备工作做好了，翻译起来才可能得心应手，势如破竹，才可能做到准确、通顺和流畅。

力求直译，字字落实。

所谓“直译”，就是在翻译的时候，跟原文一句话对着一句话，甚至一个字对着一个字的翻译。对原文的内容不能随意增减，不能笼统地说一下原文的大意，更不能加以分析评述，一定要使译文准确地反映原文的意思。例如：

“吴广素爱人，士卒多为用者。将尉醉，广故数言欲亡，忿尉，令辱之，以激怒其众。”（《陈涉世家》）

译文：吴广一向爱护别人，士兵中有很多愿意给他使唤的人。将尉喝醉了酒，吴广故意多次扬言要逃跑，使将尉愤怒，让他侮辱自己，用来激怒众人。

不能这样译：①吴广平时经常替士兵们做好事，关心别人的困难，他因而获得了众人的爱戴，大家都愿意为他帮忙，替他跑腿……（这些译法没有歪曲原意，但是增添了许多原文中没有的话，有的地方译得还不准确，如把“多为用”译成“都愿意为他帮忙……”不符合翻译的要求。）②吴广一向爱护别人，他这种做法得人心，好多士兵愿意为他服务……（这种译法，凭空加入分析评论，与原文陈述的语气不一致，也不符合翻译的要求。）

我们要求尽量直译，因为直译法不仅基本上能做到逐句对译，字字落实，使初学古文者易于养成掌握古今汉语对应规律的能力，而且不易漏译，即使有误译的地方也容易检查发现。

直译的缺点在于不易使句子表达得流畅，有些句子还难于直译，因为古今汉语在用词造句方面总有自己的特点，有时就无法使二者完全对应吻合。我们如果遇到无法直译或者直译之后非常别扭的句子，那就采用意译法。

意译是说根据原文的大意来翻译。但是意译并不是说可以笼统地翻译，还得依据原文的字、词、句的大意对照翻译，更不允许随意增删原文的意思，或者随意地加以分析评论。例如：

“吾入关，秋毫不敢有所近……”（《鸿门宴》）

直译：我攻进函谷关，连鸟兽在秋天初生的细毛都不敢有什么接近。

意译：当我攻进函谷关之后，连最细小的东西都不敢据为己有。

不能这样译：①我打进函谷关之后，不管是金银财宝、珍珠玛瑙，还是仓库粮食、大批绸缎，甚至连一支鸡、一条狗都不敢私自占有。（这种译法，随意增添原文中没有的内容，不合要求。）②我打进函谷关，不拿群众一针一线。（这种译法，既没有如实反映原文的内容，又违背了时代特点，拿今天人民军队的优良作风去乱套古代军事头领自我粉饰的做法，当然更不合要求。）

意译法的好处能使句意通顺，易于理解。对于初学者来说，应该尽量采取直译法，实在不好直译时才用意译法，这样做有助于对原文的词义、词序和句式的理解，有助于提高阅读文言文的能力。

3.严格遵循古代汉语的一般规律去正确理解词义、词序和句式。

我们已经掌握了古代汉语的一般规律，在翻译文言文时都要具体灵活地拿来应用。例如：根据古今词义的变化，古字通假和词类活用等规律，对于每一个词要确切了解其在具体语言环境中的意义，切忌望文生义；根据古代汉语一词多义的规律，一个词的确切的特定的意义必须放在整个句子里，联系上下文的文意来推求，切忌孤立地进行判断；根据古代汉语的判断句式、被动句式、定语后置、宾语前置等特殊句式或词序的规律，对应地译为现代汉语，而不能把它们置之度外，随心所欲。由于古代汉语中单音节词多，在翻译时，应补充为相应的现代汉语的双音节词或多音节词。

试以《孟子·告子上》中孟子所说的一段话为例，看看该怎样翻译：

“今有无名之指，屈而不信，非疾痛害事也。如有能信之者，则不远秦楚之路，为指之不若人也。指不若人，则知恶之；心不若人，则不知恶；此之谓不知类也。”

[译文]假设有个人，他的无名指弯曲了，不能伸直，既不很疼，又不碍事。如果有能使它伸直的人，就一定不把到秦国或者到楚国的路程看得遥远（而去医治），因为一个手指不如别人呀。一个手指不如别人，便知道厌恶它；自己思想不如别人，却不知道厌恶它；不分轻重主次说的就是这类的情况啊。

这段话中的“今”，在古代汉语中如果不指实事而指设喻的话，可以用“假设、如果”来翻译，而不必作“如今”讲。文中的“信”字，要按一般的词义“相信”讲显然是不通的，这就该想到“欲信大义于天下”（《隆中对》）这个常见通假字“伸”字了。文中的“疾”——一般当“疾病”讲，但是无名指屈而不伸总还是一种病态吧，怎能说“非疾”呢？联系上下文意，就可判断这儿用“疾”这个多义词的另一义项“猛烈”的意思了，所谓“疾风知劲草”、“声非加疾也，而闻者彰”（《劝学》），都是用的此意。文中的“远”，是形容词用作意动词，这稍加辨析就知道。文中的“为”，当然不能作动词“成为”、“变作”或者介词“被”、“替”等来讲，这只要贯串上下文意，就能判定它相当于现代的“由于、因为”，它所构成的介词结构表示动作的原因。至于两个“则”字，前一个是顺接两件紧系的事，相当于现代的“就、便”，后一个是反接两件紧系的事，相当于现代的“可是、却”。而文中两个“恶”字，当然不能作“罪恶”、“恶劣”讲，因为“恶之”的“恶”带宾语“之”，显然是该作动词“厌恶（wù）”讲，才能使文意贯通。最后一句话“此之谓不知类也”，是古代汉语中常见的一种固定格式，“此”是“谓”的宾语，中间加“之”作为宾语前置的标志。“此之谓也”等于说“说的就是这个啊”，“夫子之谓也”等于说“说的就是先生您啊”，“我之谓也”等于说“说的就是我啊”。在本文这段话中，实际是说“不知类此之谓也”，而为了强调“此之谓”，整个句子又作了主谓的倒置。“不知类”就是“不懂事”，可以根据上文语意，意译为“不分轻重主次”。

翻译时，省略的句子成分一般要补出。

我们知道，在一定的条件下，文言的句子往往省略某些句子成分，有主语的省略、动词或形容词的谓语的省略、宾语的省略、介词的省略、介词宾语的省略等。凡是由于句子成分的省略而影响文意的表达或者容易使人误解时，在译文中一律应将本应当有的被省略的句子成分加以补足。例如：

①（ ）（ ）微指左公处，（ ）

（ ）则席地倚墙而坐……（《左忠毅公逸事》）

（应译为：禁卒稍稍地指一下左公所在的地方，左公正用泥地当席子，背靠着墙壁坐在那儿……）

②婉贞于是率诸少年结束而出（ ），（ ）

皆（ ）玄衣（ ）白刃。（《冯婉贞》）

（应译为：婉贞因此带领一群青年人整好装束而出村，他们都穿着黑色的衣服，拿着雪亮的钢刀。）

③每字为一印，火烧（ ）令（ ）坚。（《活板》） （应译为：每个字做成一个字印，用火烧它，使它坚固。）

由于文言中往往不用量词，翻译时也应补译出来，如句③。

要按现代汉语的语法习惯调整词序。

文言中有不少特殊的词序和句式，其中特别是动宾倒置、介宾倒置、介词结构后置、定语后置、主谓倒置，在译成现代汉语时要将语序调整，以适合现代汉语的语法习惯。例如：

①彼不我恩也。（他不好好待我。） （《童区寄传》）

②何以效之？（拿什么验证它？） （《订鬼》）

③具告以事。（把事情全部告诉了他。） （《鸿门宴》）

④村中少年好事者驯养一虫。（村子里有个喜好多事的年轻人驯养着一只蟋蟀。） （《促织》）

⑤军书十二卷，卷卷有爷名。（很多卷的征兵名册，每卷都有阿爷的名字。） （《木兰诗》）

⑥寒暑易节，始一反焉。（冬夏变换了季节，才回来一趟。） （《愚公移山》）

除了上述几点之外，还应考虑到以下几种不必翻译的情况：

遇到人名、地名、官职、朝代、年号、书名等专有名词，可以照录，不必翻译。例如：

①沛公左司马曹无伤使人言于项羽曰：“沛公欲王关中，使子婴为相，珍宝尽有之。”（《鸿门宴》）

[译文]沛公左司马曹无伤派人对项羽说道：“沛公想在关中地区称王，让子婴做宰相，把珍器宝物完全占为己有。”

②康熙五十一年三月，余在刑部狱，见死而由窦出者日三四人。（《狱中杂记》）

[译文]康熙五十一年三月，我被关在刑部的监狱，每天看到死掉而从墙洞中拖出去的有三四个人。

遇到一些表敬或谦称的词，可以不译。例如：
①太后曰：“敬诺。” （太后说：“好的。”） （《触龙说赵太后》）
②窃为大王不取也。（认为大王不该这样做。） （《鸿门宴》）

遇到在句里只起语法作用而无实在意义的文言虚词，可以不译。例如：
①夫六国与秦皆诸侯，其势弱于秦……（六国与秦国都是诸侯国，但六国的实力比秦国弱。） （《六国论》）
②一夫不耕，或受之饥。（一个男人不耕作，有人就挨饿。） （《论贵粟疏》）
③珍宝尽有之。（把珍器宝物完全占有）（《鸿门宴》）
④谁得而族灭也？（谁能够族灭它呢？）（《阿房宫赋》）
⑤则吾恐秦人食之不得下咽也。（那么我恐怕秦国人吃饭也不能朝下咽了。） （《六国论》）

遇到在现代汉语中还经常应用的成语和习惯用语，一般可以不译，因为不译不仅不妨碍人们对这些词语的理解，而且可以使译文通顺，保持原文的语言特色。例如“气象万千”、“千钧一发”、“叱咤风云”、“走投无路”等成语，“于是”、“因而”、“其余”、“足以”等连用的习惯用语，就是如此。
\newpage